\section{Introduction}
\label{sec:tpc37-introduction}

The preceding papers reduce one calibrated prime--M\"obius residual to a
specific orbit--sliced energy problem.  TPC-25 shows, in an earlier
one-sided shell, that physical coefficient provenance must be retained
before taking absolute values \citep{WangTPC25}.  TPC-35 identifies the
complete product--slope spectrum and the critical three-modulus alias
scale \citep{WangTPC35}.  TPC-36 then proves that the actual static row
mask has only $X^{o(1)}$ signed projective cost, and relocates the
remaining difficulty to the weighted coherent functional
\begin{equation}
 \frac1{q-1}\sum_\chi
 |\widehat b(\bar\chi)|^2
 \left|\sum_s w(s)\widehat R_s(\chi)\right|^2
 \label{eq:tpc37-intro-inherited-functional}
\end{equation}
after incomplete ranges and integer aliases are restored
\citep{WangTPC36}.  The present paper separates three coordinate effects
that are entangled in \eqref{eq:tpc37-intro-inherited-functional}:
physical zero residues, the nonconstant complementary-factor weight, and
the proper faces of a three-modulus completion.

The inherited sufficient estimate is
\begin{equation}
 \boxed{
 V_L+V_R\ll X^{o(1)}\frac{Q^3}{J}.}
 \label{eq:tpc37-intro-orbit-gate}
\end{equation}
The sign-blind exact-equality energy has the larger baseline
\begin{equation}
 V_{\rm eq}^{\rm abs}\ll X^{o(1)}Q^3J.
 \label{eq:tpc37-intro-absolute-baseline}
\end{equation}
At the high--\(\beta\) endpoint,
\begin{equation}
 Q=X^{267/400+o(1)},\qquad
 J=X^{133/400+o(1)},\qquad
 T=X^{193/500+o(1)},
 \label{eq:tpc37-intro-scales}
\end{equation}
so the available ratio between
\eqref{eq:tpc37-intro-orbit-gate} and
\eqref{eq:tpc37-intro-absolute-baseline} is exactly $J^{-2}$ in
energy.  Equivalently, the allowance above the previously isolated
Gram-copy diagonal $X^{o(1)}Q^2J$ is
\begin{equation}
 \frac{Q^3/J}{Q^2J}
 =\frac Q{J^2}=X^{1/400+o(1)}.
 \label{eq:tpc37-intro-margin}
\end{equation}
There is no room for an untracked fixed power of $X$.

\subsection{Zero coordinates are removed on the physical equality}

Choose a prime $q\asymp J$ larger than the positive orbit support and
between the physical $d$- and $\ell$-scales.  On the part of the
actual packet satisfying
\begin{equation}
 q\mid m_\alpha j+h,
 \label{eq:tpc37-intro-singular-target}
\end{equation}
a fixed orbit time selects only $O(1+Q/q)$ row labels, whereas a fixed
row label occurs at no more than one orbit time.  This elementary
two-sided incidence, combined with the inherited row $\ell^2$-mass,
gives
\begin{equation}
 \boxed{
 V_{L,\rm sing(q)}+V_{R,\rm sing(q)}
 \ll X^{o(1)}(1+Q/q)Q^2
 \ll X^{o(1)}\frac{Q^3}{J}.}
 \label{eq:tpc37-intro-zero-face-bound}
\end{equation}
If the ultra divisor is opened as
\begin{equation}
 su=m_\alpha j+h,
 \label{eq:tpc37-intro-physical-equality}
\end{equation}
then, because $q$ is prime, the event in
\eqref{eq:tpc37-intro-singular-target} is precisely the union of the
physical faces $s=0\pmod q$ and $u=0\pmod q$.  Thus the zero
coordinates are absorbed before replacing equality by congruence.  The
argument extends to any fixed number of orbit-scale primes.

This result does not discard the terminal factorization.  On the regular
support $q\nmid m_\alpha j+h$, the point $s=1$,
$u=m_\alpha j+h$ remains present and both coordinates are units.

\subsection{Punctured centering isolates the complete trace}

After physical zero-face excision, the raw ultra residue vector $G$
satisfies $G(0)=0$.  Ordinary additive centering would generally
create a nonzero value at the excluded residue.  We instead define
\begin{equation}
 \beta=\frac1{q-1}\sum_{U\ne0}G(U),\qquad
 g(0)=0,\qquad g(U)=G(U)-\beta\quad(U\ne0).
 \label{eq:tpc37-intro-punctured-centering}
\end{equation}
If $f=F-q^{-1}\sum_D F(D)$, then for $a,s\ne0$ the raw affine
correlation satisfies the exact identity
\begin{equation}
 \mathcal C_{a,s}(F,G)
 =\mathcal C_{a,s}(f,g)
 -\beta f(-h/a)+(q-1)\overline F\,\beta.
 \label{eq:tpc37-intro-punctured-identity}
\end{equation}
Every correction is explicit and rank one in the complementary
\(s\)-coordinate.  The punctured-centered core
has the exact complete-cycle cancellation
\begin{equation}
 \sum_{s\in\mathbb F_q^\times}R_s(a)=0.
 \label{eq:tpc37-intro-complete-cycle}
\end{equation}
For the actual dyadic ultra coefficient, a standard one-M\"obius
estimate gives, for every fixed $A>0$,
\begin{equation}
 |\beta|\ll_A\frac Uq(\log X)^{-A}.
 \label{eq:tpc37-intro-beta-bound}
\end{equation}
Hence all complete residue cycles in the complementary interval vanish
from the centered core; the incomplete boundary, the displayed mean
corrections, and alias weights remain.  The logarithmic saving in
\eqref{eq:tpc37-intro-beta-bound} is not promoted here to a fixed power
saving.

\subsection{The two-character layer is a Jacobi near-isometry}

Let $w$ be the nonconstant complementary-factor weight and let
$c=w*_\times g$ on $\mathbb F_q^\times$.  The coherent transform can
be written as
\begin{equation}
 C_\chi=\sum_{x\in\mathbb F_q^\times}
 c(x)\overline{\chi(x-h)}.
 \label{eq:tpc37-intro-shifted-Mellin}
\end{equation}
Taking a second Mellin transform produces the nonprincipal Jacobi matrix
\begin{equation}
 K_{\chi,\psi}=J(\psi,\bar\chi),
 \qquad \chi,\psi\ne\mathbf1.
 \label{eq:tpc37-intro-Jacobi-matrix}
\end{equation}
We compute its Gram matrix exactly:
\begin{equation}
 \boxed{
 K^*K=KK^*=(q-1)^2I-q\,\mathbf1\mathbf1^*.}
 \label{eq:tpc37-intro-Jacobi-Gram}
\end{equation}
Thus its singular values are $q-1$, with multiplicity $q-3$, and
$1$, with multiplicity one.  The Jacobi layer is therefore a scaled
near-isometry rather than a source of free square-root cancellation.
The same calculation gives the exact energy identity
\begin{align}
 \sum_{\chi\ne\mathbf1}|C_\chi|^2
 ={}&\sum_{\psi\ne\mathbf1}
 |\widehat w(\psi)\widehat g(\psi)|^2 \notag\\
 &-\frac q{(q-1)^2}
 \left|\sum_{\psi\ne\mathbf1}
 \widehat w(\psi)\widehat g(\psi)\psi(h)\right|^2.
 \label{eq:tpc37-intro-Jacobi-energy}
\end{align}

When $w$ is the residue-count weight of an interval,
P\'olya--Vinogradov bounds its nonprincipal Mellin coefficients and saves
one full factor $q$ relative to uniform Cauchy in the unweighted
character average.  That conclusion is deliberately not transferred to
the outer physical coefficient direction: a punctured single-character
example still attains the coherent Cauchy factor.  For the nonprincipal
outer spectrum, a sufficient hypothesis is a Mellin-flatness estimate for
the physical convolution \(b*_\times f\), or an equivalent bound for the
resulting M\"obius Hadamard energy.  The principal outer boundary remains
an explicit separate term.  Neither that boundary estimate nor the
nonprincipal hypothesis is proved.

\subsection{Three-modulus proper faces localize, but the full face remains}

For three pairwise coprime primes $q_i\asymp J$, we decompose the CRT
delta into its eight principal/nonprincipal coordinate faces.  They are
orthogonal over one CRT period, with squared masses
\begin{equation}
 J^{-3},\qquad J^{-2},\qquad J^{-1},\qquad 1
 \label{eq:tpc37-intro-face-scales}
\end{equation}
according to the number of nonprincipal coordinates.  More strongly,
the coherent sum $P$ of all seven proper faces satisfies
\begin{equation}
 \|P\|_\infty\ll J^{-1},\qquad
 \|P\|_2^2\asymp J^{-1},\qquad
 \|P\|_1\asymp1.
 \label{eq:tpc37-intro-proper-norms}
\end{equation}
Grouping first by the integer determinant $F=su-N$ removes an apparent
incomplete-box boundary and preserves these exact period norms up to the
critical multiplicity
\begin{equation}
 1+\frac X{q_1q_2q_3}
 =X^{1/400+o(1)}.
 \label{eq:tpc37-intro-critical-periods}
\end{equation}
This is a fiberwise statement.  It does not commute automatically with
the coherent row and divisor sums in the physical energy.

On the genuine equality $F=0$, the proper union has amplitude
$O(J^{-1})$.  Therefore its contribution closes directly from the
sign-blind baseline:
\begin{equation}
 V_{{\rm eq},{\rm proper}}
 \ll X^{o(1)}J^{-2}Q^3J
 =X^{o(1)}\frac{Q^3}{J}.
 \label{eq:tpc37-intro-proper-equality-closure}
\end{equation}
The full three-coordinate face has amplitude $1-O(J^{-1})$ on the same
equality and retains the terminal M\"obius layer.  An artificial bounded
coefficient can also align with the $\ell^1$, rather than the
$\ell^2$, size of the proper face.  Consequently the periodic quadratic
calculation is not a coefficient-free proof of
\eqref{eq:tpc37-intro-orbit-gate}.

\subsection{Claim boundary and organization}

The unconditional advances are the physical absorption
\eqref{eq:tpc37-intro-zero-face-bound}, the exact punctured-centering and
Jacobi identities, the orthogonal CRT face calculus, determinant-first
periodization, and proper-face closure on the exact equality.  The
remaining alternatives are either separate absorption of the exposed
punctured-centering mean corrections and principal outer boundary
together with an actual nonprincipal Mellin-overlap theorem, or a
target-coupled, coefficient-weighted shifted-fiber Bessel estimate for
the unsplit raw physical output.  We establish neither route.
In particular, this paper proves no estimate for the full
regular coherent face, no affine Chowla theorem, no parity breach, no
positivity or prime-pair main term, no Hardy--Littlewood asymptotic, and
no twin-prime theorem.

\Cref{sec:tpc37-inherited-interface} fixes the physical object and the
normalizations.  \Cref{sec:tpc37-physical-zero-faces} removes the physical
zero coordinates.  \Cref{sec:tpc37-punctured-centering} proves the
punctured identities and complete-cycle cancellation.
\Cref{sec:tpc37-two-character} gives the exact weighted Jacobi spectrum.
\Cref{sec:tpc37-three-modulus} develops the three-modulus face calculus
and determinant-first reassembly.  The final sections state the remaining
gate, certificate scope, and formal nonclaims.
